\subsection{Solutions des ateliers}
\begin{enumerate}
\item Les images de $0,1,\ii$ sont $- 2 \mathrm{i}, 1 - \mathrm{i}, -1 - \mathrm{i}$. Pour tout couple, $T(z)-T(w)=(1+\ii)(z-w)$ ; son module vaut $|1+\ii||z-w|=\sqrt2|z-w|$. La translation n'intervient pas dans une différence.
\item Résoudre l'égalité affine donne $z=(u+2\ii)/(1+\ii)$. Pour $u=3+\ii$, on obtient $z=3$ et la substitution redonne $3+\ii$. Le coefficient $1+\ii$ étant non nul, la division détermine exactement une entrée.
\item En élevant les deux distances positives au carré, on obtient $(x+2)^2+y^2=x^2+(y-2)^2$, soit $x+y=0$. C'est la médiatrice du segment d'extrémités $(-2,0)$ et $(0,2)$. Son milieu $(-1,1)$ satisfait bien l'équation.
\end{enumerate}
